\chapter{Preliminaries for Large Sample Theory}
\label{cha:preliminaries_large_sample}
We review basic facts about stochastic convergence and limit theorems
that we will draw on in our later discussions of large-sample theory for
statistical methodology. For additional reading, consider the initial
chapters of Ferguson (1996) and van der Vaart (1998).

\section{Stochastic Convergence}

Let $X_n$ be a sequence of $k$--dimensional random vectors on a probability space $(\Omega, \mathcal{A}, \mathsf{P})$. We distinguish between different ``modes" of how this sequence approaches a limit $X$.

\subsection*{1. Almost Sure Convergence ($\xrightarrow{a.s.}$)}
This is the strongest form of pointwise convergence. The sequence converges for ``all" outcomes $\omega$, except possibly for a set of probability zero.

\[ X_n \xrightarrow{a.s.} X \iff \mathsf{P}\left( \{ \omega : \lim_{n \to \infty} X_n(\omega) = X(\omega) \} \right) = 1 \]

\textit{Characterization via Supremum:}
\[ X_n \xrightarrow{a.s.} X \iff \forall \epsilon > 0 : \ \mathsf{P}\left(\sup_{m \geq n} \|X_m - X\| > \epsilon\right) \to 0 \text{ as } n \to \infty. \]

\subsection*{2. Convergence in $r$--th Mean ($\xrightarrow{r}$)}
Often used with $r=2$ (Quadratic Mean). It measures convergence using the expected value of the distance raised to the power $r$.

\[ X_n \xrightarrow{r} X \iff \mathsf{E}[\|X_n - X\|^r] \to 0 \quad \text{as} \quad n \to \infty. \]
\begin{itemize}
    \item Requires $\mathsf{E}[\|X_n\|^r] < \infty$ and $\mathsf{E}[\|X\|^r] < \infty$.
    \item \textbf{Implication:} If $X_n \xrightarrow{r} X$, then $X_n \xrightarrow{r'} X$ for any $0 < r' \le r$.
\end{itemize}

\subsection*{3. Convergence in Probability ($\xrightarrow{p}$)}
This is the standard mode for statistical consistency. It demands that the probability of the estimator $X_n$ deviating from $X$ by any fixed amount $\epsilon$ goes to zero.

\[ X_n \xrightarrow{p} X \iff \forall \epsilon > 0 : \mathsf{P}(\|X_n - X\| > \epsilon) \to 0 \text{ as } n \to \infty. \]

\subsection*{4. Convergence in Distribution ($\xrightarrow{d}$)}
This is the weakest form. It does not require $X_n$ and $X$ to be close to each other value-wise; it only requires their \textit{distribution functions} (CDFs) to become similar.

\[ X_n \xrightarrow{d} X \iff F_n(x) \to F(x) \quad \forall x \text{ where } F \text{ is continuous.} \]
\begin{note}[]
This is the mode of convergence used in the Central Limit Theorem.
\end{note}

\subsection*{Relationships and Properties}

\subsubsection{Implications}
\begin{itemize}
    \item \textbf{Top Tier:} $\xrightarrow{a.s.}$ and $\xrightarrow{r}$ are the strongest. They do not imply each other directly (without extra conditions).
    \item \textbf{Middle Tier:} Both $\xrightarrow{a.s.}$ and $\xrightarrow{r}$ imply $\xrightarrow{p}$.
    \item \textbf{Bottom Tier:} $\xrightarrow{p}$ implies $\xrightarrow{d}$.
\end{itemize}

\subsection*{Useful Tools}

\textbf{1. Subsequence Characterization}
$X_n \xrightarrow{p} X$ if and only if every subsequence $(n_k)$ contains a further sub-subsequence $(n_{k_\ell})$ such that $X_{n_{k_\ell}} \xrightarrow{a.s.} X$.
\begin{quote}
    \textit{Usage:} If you find proving $\xrightarrow{p}$ directly is hard, but you can prove $\xrightarrow{a.s.}$ for subsequences, you can use this equivalence.
\end{quote}

\textbf{2. Convergence to a Constant}
If the limit is a constant $c$ (non-random), then convergence in distribution is equivalent to convergence in probability:
\[ X_n \xrightarrow{d} c \iff X_n \xrightarrow{p} c \]

\textbf{3. The Cramér-Wold Device (for $k > 1$)}
Used to prove convergence in distribution for random \textbf{vectors} by reducing the problem to scalar (1D) linear combinations.

\[ X_n \xrightarrow{d} X \iff a^\top X_n \xrightarrow{d} a^\top X \quad \forall a \in \mathbb{R}^k \]
\begin{quote}
    \textit{Why it works:} It relies on Characteristic Functions (Fourier transforms of distributions). If the characteristic functions match for all linear combinations $t(a^\top X)$, the joint distributions must match.
\end{quote}

\section{Limit Theorems}

We consider an i.i.d. sequence $(X_i)_{i=1}^{\infty}$ with finite mean $\mu = \mathsf{E}[X_i]$. Let $\overline{X}_n$ be the sample mean.

\subsection*{1. Law of Large Numbers (LLN)}
The LLN guarantees that sample averages converge to the population mean.

\begin{itemize}
    \item \textbf{Strong LLN:} $\overline{X}_n \xrightarrow{a.s.} \mu$. (Convergence with probability 1).
    \item \textbf{Weak LLN:} $\overline{X}_n \xrightarrow{p} \mu$. (Convergence in probability).
\end{itemize}
\begin{remark}[]
Since $\xrightarrow{a.s.}$ implies $\xrightarrow{p}$, the Strong LLN implies the Weak LLN.
\end{remark}

\subsection*{2. Central Limit Theorem (CLT)}
The CLT describes the distribution of the error $\overline{X}_n - \mu$.

\textbf{Univariate ($k=1$):}
If $\sigma^2 = \mathsf{Var}[X_i] \in (0, \infty)$, then:
\[ \sqrt{n}\left(\frac{\overline{X}_n - \mu}{\sigma}\right) \stackrel{d}{\longrightarrow} \mathcal{N}(0, 1). \]

\textbf{Multivariate ($k > 1$):}
If the covariance matrix $\Sigma = \mathsf{Var}[X_i]$ exists, then:
\[ \sqrt{n}(\overline{X}_n - \mu) \stackrel{d}{\longrightarrow} \mathcal{N}_k(0, \Sigma). \]

\subsection*{3. CLTs for Triangular Arrays}
Used when the data generating distribution changes with $n$.
Let $X_{n,1}, \dots, X_{n,n}$ be independent random variables for each row $n$, with:
\[ \mu_{n,i} = \mathsf{E}[X_{n,i}] \quad \text{and} \quad \sigma_{n,i}^2 = \mathsf{Var}[X_{n,i}]. \]
Let $\sigma_n^2 = \sum_{i=1}^n \sigma_{n,i}^2$ be the total variance of the sum in row $n$.

\textbf{Lindeberg-Feller CLT:}
The standardized sum converges to a standard normal:
\[ \frac{1}{\sigma_n}\sum_{i=1}^n (X_{n,i} - \mu_{n,i}) \stackrel{d}{\longrightarrow} \mathcal{N}(0,1) \]
\textit{Provided the \textbf{Lindeberg Condition} holds:}
\[ \forall \epsilon > 0 : \frac{1}{\sigma_n^2} \sum_{i=1}^n \mathsf{E}\left[ (X_{n,i} - \mu_{n,i})^2 \cdot \bs{1}_{\{|X_{n,i} - \mu_{n,i}| > \sigma_n \epsilon\}} \right] \underset{n \to \infty}{\longrightarrow} 0. \]
\textit{Intuition:} The contribution of the ``tails" (extreme values) to the total variance must be negligible.

\textbf{Lyapunov CLT (Simpler Condition):}
The Lindeberg condition is satisfied (and thus the CLT holds) if the \textbf{Lyapunov Condition} holds:
\[ \frac{1}{\sigma_n^3} \sum_{i=1}^n \mathsf{E}[|X_{n,i} - \mu_{n,i}|^3] \underset{n \to \infty}{\longrightarrow} 0. \]
\textit{Usage:} It is often easier to calculate the 3rd moments than to evaluate the indicator function in the Lindeberg condition.

%\section{Tools for Convergence: Skorokhod and Continuous Mapping}

\subsection*{The Skorokhod Representation Theorem}

The \textbf{Skorokhod Representation Theorem} is a powerful theoretical tool. While convergence in distribution ($X_n \xrightarrow{d} X$) is the ``weakest" form of convergence (implying only that CDFs converge), this theorem allows us to ``upgrade" this convergence to almost sure convergence ($\xrightarrow{a.s.}$) by constructing new random variables on a different probability space.

%\textbf{Lecturer's Intuition:}
Think of this as a ``technical device" or a bridge. It allows us to prove properties about weak convergence ($\xrightarrow{d}$) using the stronger machinery of almost sure convergence ($\xrightarrow{a.s.}$). It essentially says: ``If the distributions converge, I can simulate these random variables in a way that the values actually converge point-wise."

\begin{theo}[Skorokhod Representation]
Suppose $X_n \xrightarrow{d} X$. Then there exist random vectors $\tilde{X}_n$ and $\tilde{X}$ defined on a common probability space $(\tilde{\Omega}, \tilde{\mathcal{A}}, \tilde{P})$ such that:
\begin{enumerate}
    \item \textbf{Distributions Match:} $\tilde{X}_n \stackrel{d}{=} X_n$ for all $n$, and $\tilde{X} \stackrel{d}{=} X$.
    \item \textbf{Strong Convergence:} $\tilde{X}_n \xrightarrow{a.s.} \tilde{X}$ as $n \to \infty$.
\end{enumerate}
\end{theo}

\begin{proof}[The Quantile Transform]
For the \(1\)--dimensional case ($k=1$), the proof is constructive and relies on the \textbf{Probability Integral Transform}.
\begin{enumerate}
    \item Take a single random source $U \sim \text{Unif}(0, 1)$.
    \item Define the random variables using the quantile functions (generalized inverses) of the CDFs:
    \[ \tilde{X}_n := F_{X_n}^{-1}(U) \quad \text{and} \quad \tilde{X} := F_X^{-1}(U). \]
    \item Since $X_n \xrightarrow{d} X$, we know $F_{X_n}(x) \to F_X(x)$ at continuity points. This implies the quantile functions converge point-wise.
    \item Because $\tilde{X}_n$ and $\tilde{X}$ are built from the \textit{exact same} $U$, they are highly dependent (coupled), which forces them to converge almost surely.
\end{enumerate}

\textbf{Warning:} We lose all information about the joint relationships (e.g., independence) of the original sequence $X_n$. Skorokhod preserves the \textit{marginal} distributions of $X_n$ and $X$, but creates a specific dependence structure to force path-wise convergence.

For $k > 1$, the construction is significantly more complex; see Billingsley (1986).
\end{proof}

\subsection*{The Continuous Mapping Theorem (Mann--Wald)}

This theorem answers a fundamental question: 
\begin{quote}
If $X_n$ converges to $X$, does $g(X_n)$ converge to $g(X)$?
\end{quote}

This is the stochastic analog of the definition of continuity. For deterministic numbers, continuity means $x_n \to x \implies g(x_n) \to g(x)$. This theorem states that this preservation of convergence holds for our three stochastic modes.

\begin{theo}[Continuous Mapping Theorem / Mann--Wald]

\label{thm:cmt}

Let $g: \mathbb{R}^k \to \mathbb{R}^m$ be a function that is continuous almost everywhere with respect to the distribution of $X$ (i.e., $\mathsf{P}(X \in D_g) = 0$, where $D_g$ is the set of discontinuity points of $g$). Then:
\begin{enumerate}
    \item $X_n \xrightarrow{a.s.} X \implies g(X_n) \xrightarrow{a.s.} g(X)$
    \item $X_n \xrightarrow{p} X \implies g(X_n) \xrightarrow{p} g(X)$
    \item $X_n \xrightarrow{d} X \implies g(X_n) \xrightarrow{d} g(X)$
\end{enumerate}
\end{theo}



\begin{proof}[]
    \leavevmode
\begin{enumerate}[]
    \item 
\textbf{Almost Sure Convergence ($\xrightarrow{a.s.}$):}
There is nothing to do.
\begin{itemize}
    \item By definition, there exists a set $A$ with $\mathsf{P}(A)=1$ where $X_n(\omega) \to X(\omega)$ pointwise.
    \item By the definition of continuity, for every $\omega \in A$, $g(X_n(\omega)) \to g(X(\omega))$.
\end{itemize}
    Thus, $g(X_n) \xrightarrow{a.s.} g(X)$.
\item 
\textbf{Convergence in Probability ($\xrightarrow{p}$):}
Use the Subsequence Characterization.
\begin{itemize}
    \item Instead of struggling with $\epsilon$-$\delta$ definitions for probability, recall: $Y_n \xrightarrow{p} Y$ if and only if every subsequence has a further sub-subsequence that converges almost surely.
    \item Take any subsequence of $X_n$. It has a sub-subsequence converging a.s.
    \item By result (1), applying $g$ preserves this a.s. convergence.
\end{itemize}
    Therefore, $g(X_n) \xrightarrow{p} g(X)$.

    \item 
\textbf{Convergence in Distribution ($\xrightarrow{d}$):}
Use Skorokhod as a bridge.
\begin{itemize}
    \item Proving this directly with integrals and CDFs is messy.
    \item Instead, apply the \textbf{Skorokhod Representation Theorem}.
    \item Switch from $X_n \xrightarrow{d} X$ to the coupled version $\tilde{X}_n \xrightarrow{a.s.} \tilde{X}$.
    \item Apply result (1) (continuous mapping for a.s.) to get $g(\tilde{X}_n) \xrightarrow{a.s.} g(\tilde{X})$.
    \item Since almost sure convergence implies convergence in distribution, we have $g(\tilde{X}_n) \xrightarrow{d} g(\tilde{X})$.
    \item Since distributions are preserved ($\tilde{X}_n \stackrel{d}{=} X_n$), we conclude $g(X_n) \xrightarrow{d} g(X)$.
\end{itemize}
\end{enumerate}    
\end{proof}


\subsection*{Slutsky's Theorem}

Slutsky's theorem is a workhorse in asymptotic statistics. It allows us to perform algebraic manipulations on sequences of random variables (like addition and multiplication) while preserving their convergence properties, \textit{provided} that one of the sequences converges to a constant.

\begin{theo}[Slutsky's Theorem]
    \label{thm:slutsky}
Let $X_n \xrightarrow{d} X$.
\begin{enumerate}[label=\alph*)]
    \item \textbf{Univariate case ($k=1$):} If $A_n \xrightarrow{p} a$ and $B_n \xrightarrow{p} b$ for constants $a, b \in \mathbb{R}$, then the affine transformation converges:
    \[ A_n X_n + B_n \xrightarrow{d} aX + b. \]
    \item \textbf{General case ($k \ge 1$):} If $X_n \xrightarrow{d} X$ and $Y_n \xrightarrow{p} c$ for a constant vector $c$, then the joint vector converges in distribution:
    \[ \begin{pmatrix} X_n \\ Y_n \end{pmatrix} \xrightarrow{d} \begin{pmatrix} X \\ c \end{pmatrix}. \]
    Consequently, by the Continuous Mapping Theorem, for any continuous function $g$:
    \[ g(X_n, Y_n) \xrightarrow{d} g(X, c). \]
\end{enumerate}
\end{theo}

\subsubsection*{The ``Stacking" Argument}
Why is the condition $Y_n \xrightarrow{p} c$ (constant) so important?
\begin{itemize}
    \item \textbf{The Problem:} generally, if $X_n \xrightarrow{d} X$ and $Y_n \xrightarrow{d} Y$ (random), we know \textbf{nothing} about the convergence of the pair $(X_n, Y_n)$. They could be independent, perfectly correlated, or have a dependency structure that changes with $n$. Without knowing their joint dependence, we cannot determine the limit of $g(X_n, Y_n)$.
    \item \textbf{The Solution:} If $Y_n$ converges to a \textit{constant} $c$, the dependence ``decouples" in the limit. The ``stacking" of a random vector $X_n$ and a vector converging to a constant always converges jointly. This allows us to treat $Y_n$ as if it were the constant $c$ for the purposes of asymptotic algebra.
\end{itemize}

\begin{quote}
    \textbf{Practical use in Statistics:} We often derive a limit theorem involving a true parameter (e.g., standardizing by $\sigma$). In practice, we must replace $\sigma$ with an estimator $S_n$. Slutsky's theorem guarantees that as long as $S_n$ is \textbf{consistent} ($S_n \xrightarrow{p} \sigma$), the asymptotic distribution of our test statistic remains unchanged.
\end{quote}


\begin{ex}[Asymptotic Normality of the \(t\)--Statistic]

This is the classic application of Slutsky's theorem. We want to show that replacing the true standard deviation $\sigma$ with the sample standard deviation $S_n$ in the CLT does not change the limit distribution.

\textbf{Setup:} Let $X_1, X_2, \dots$ be i.i.d. with $\mathsf{E}[X_i] = \mu$ and $\mathsf{Var}[X_i] = \sigma^2 \in (0, \infty)$.

\subsubsection*{Consistency of Sample Moments (LLN)}
By the Weak Law of Large Numbers (WLLN):
\[ \overline{X}_n \xrightarrow{p} \mu \quad \text{and} \quad \overline{X^2_n} = \frac{1}{n} \sum_{i=1}^n X_i^2 \xrightarrow{p} \mathsf{E}[X_i^2] = \mu^2 + \sigma^2. \]
Since both converge in probability to constants, their joint vector converges:
\[ \left(\overline{X}_n, \overline{X^2_n}\right) \xrightarrow{p} (\mu, \mu^2 + \sigma^2). \]

\subsubsection*{Consistency of Variance Estimator (Continuous Mapping)}
We express the sample variance using the identity:
\[ \frac{1}{n} \sum_{i=1}^{n} (X_i - \overline{X}_n)^2 = \overline{X^2_n} - (\overline{X}_n)^2. \]
Define the function $g(u, v) = v - u^2$, which is continuous. Applying the Continuous Mapping Theorem to the result from Step 1:
\[ \overline{X^2_n} - (\overline{X}_n)^2 \xrightarrow{p} (\mu^2 + \sigma^2) - (\mu)^2 = \sigma^2. \]
\begin{remark}[on Bias Correction]
    
We typically use $S^2 = \frac{1}{n-1} \sum (X_i - \overline{X})^2$. Since $\frac{n}{n-1} \to 1$, the factor is asymptotically negligible. Thus:
\[ S_n^2 \xrightarrow{p} \sigma^2 \implies S_n \xrightarrow{p} \sigma. \]
This formally proves that $S_n$ is a consistent estimator for $\sigma$.
\end{remark}

\subsubsection*{The \(t\)--Statistic (CLT + Slutsky)}
We want the limit of the t-statistic: $T_n = \frac{\sqrt{n}(\overline{X}_n - \mu)}{S_n}$.
We can rewrite this as a product:
\[ \frac{\sqrt{n}(\overline{X}_n - \mu)}{S_n} = \underbrace{\frac{\sqrt{n}(\overline{X}_n - \mu)}{\sigma}}_{\xrightarrow{d} \mathcal{N}(0,1) \text{ by CLT}} \cdot \underbrace{\frac{\sigma}{S_n}}_{\xrightarrow{p} 1 \text{ by CMT}}. \]

\begin{itemize}
    \item The first term converges in distribution to $\mathcal{N}(0,1)$ by the Central Limit Theorem.
    \item The second term converges in probability to $1$ because $S_n \xrightarrow{p} \sigma$.
\end{itemize}

Applying Slutsky's Theorem (product rule):
\[ T_n \xrightarrow{d} \mathcal{N}(0, 1) \cdot 1 = \mathcal{N}(0, 1). \]

\begin{note}
This confirms that for large samples, we can effectively ``ignore" the fact that we estimated the variance; the \(t\)--statistic behaves exactly like a standard normal variable.
\end{note}
\end{ex}

\section{The Delta Method}

While the Continuous Mapping Theorem tells us that applying a continuous function preserves convergence (i.e., $X_n \xrightarrow{d} X \implies g(X_n) \xrightarrow{d} g(X)$), it does not tell us about the \textit{rate} of convergence or the resulting distribution of the fluctuations.

\paragraph{Error Propagation}%
\label{par:Error Propagation}

Think of this as the statistical equivalent of ``error propagation" in numerical analysis.
Suppose $X_n$ is an estimator for a parameter $b$. As $n \to \infty$, the estimator becomes more accurate, meaning the error $X_n - b \to 0$. However, if we scale this error by a sequence $a_n$ (often $\sqrt{n}$), the scaled error stabilizes to a limit distribution $Z$ (e.g., a standard normal).
The Delta Method answers the question: 
\begin{quote}
If we transform our estimator via a function $g$, how do the estimation errors of $g(X_n)$ behave relative to the target $g(b)$?
\end{quote}

The method relies on a \textbf{linearization} (first-order Taylor approximation) of the function $g$. If $g$ is differentiable, linearizing it locally around the parameter $b$ allows us to transfer the central limit behavior of $X_n$ to $g(X_n)$.

\begin{theo}[The Delta Method / Cramér's Theorem]\label{thm:delta_method}
Let $(X_n)_{n\in\mathbb{N}}$ be a sequence of random vectors in $\mathbb{R}^m$. Let $a_n \in \mathbb{R}$ and $b \in \mathbb{R}^m$ be such that $a_n \to \infty$ and
\[ Z_n := a_n(X_n - b) \xrightarrow{d} Z. \]
If $g : \mathbb{R}^m \to \mathbb{R}^k$ is differentiable at $b$, with Jacobian matrix $g'(b) \in \mathbb{R}^{k \times m}$, then
\[ a_n(g(X_n) - g(b)) \xrightarrow{d} g'(b)Z. \]
\end{theo}

\begin{note}[on Dimensions]
If $g$ maps $\mathbb{R}^m \to \mathbb{R}^k$, the Jacobian $g'(b)$ is a $k \times m$ matrix. The limit is the random vector $Z$ linearly transformed by this matrix.
\end{note}

\begin{note}[The Normal Case]
The most common application is when the limit $Z$ is multivariate normal. If $Z \sim \mathcal{N}_m(0, \Sigma)$, then by the properties of linear transformations of Gaussian vectors:
\[ g'(b)Z \sim \mathcal{N}_k(0, g'(b)\Sigma g'(b)^T). \]
In this specific context, the result is often referred to as \emph{Cramér's Theorem}.
\end{note}

\begin{proof}
We consider $m=1$ and show an explicit difference quotient. The case $m>1$ is proven analogously; see van der Vaart (1998, Theorem 3.1).

\textbf{Skorokhod Representation.}
Working directly with convergence in distribution is abstract because the random variables as functions can behave erratically. To make the analysis ``pedestrian" (standard calculus), we use the Skorokhod representation theorem.
There exist $\tilde{X}_n, \tilde{Z}_n, \tilde{Z}$ on a common probability space such that:
\begin{itemize}
    \item $\tilde{Z}_n \stackrel{d}{=} Z_n$ and $\tilde{Z} \stackrel{d}{=} Z$.
    \item $\tilde{Z}_n \xrightarrow{a.s.} \tilde{Z}$.
\end{itemize}
We define $\tilde{X}_n$ via the relation $\tilde{Z}_n = a_n(\tilde{X}_n - b)$. Since $\tilde{Z}_n$ converges to a finite limit and $a_n \to \infty$, it must be that $\tilde{X}_n \xrightarrow{a.s.} b$.

\textbf{Linearization.}
We can now analyze the difference quotient directly using pathwise convergence.
\[
a_n\big(g(\tilde{X}_n)-g(b)\big) =
\begin{cases}
\frac{g(\tilde{X}_n)-g(b)}{\tilde{X}_n-b}(\tilde{X}_n-b)a_n, & \text{if } \tilde{X}_n \neq b \\
0, & \text{if } \tilde{X}_n = b
\end{cases}
\]
As $n \to \infty$:
\begin{itemize}
    \item The term $(\tilde{X}_n - b)a_n = \tilde{Z}_n \xrightarrow{a.s.} \tilde{Z}$.
    \item Since $\tilde{X}_n \to b$ and $g$ is differentiable, the difference quotient $\frac{g(\tilde{X}_n)-g(b)}{\tilde{X}_n-b} \xrightarrow{a.s.} g'(b)$.
\end{itemize}
Thus, the product converges almost surely to $g'(b) \cdot \tilde{Z}$. Since distributional convergence is preserved under this representation,
\[ a_n(g(X_n) - g(b)) \xrightarrow{d} g'(b)Z. \]
\end{proof}

\begin{ex}[Squared Sample Mean]
    
Let $X_1, \dots, X_n$ be i.i.d. with $\mathsf{E}[X_i] = \mu$ and $\mathsf{Var}[X_i] = \sigma^2 < \infty$. By the CLT, the estimation error of the mean behaves as:
\[ \sqrt{n}(\overline{X}_n - \mu) \stackrel{d}{\longrightarrow} \mathcal{N}(0, \sigma^2). \]
We wish to estimate $\mu^2$ using the estimator $(\overline{X}_n)^2$.

\paragraph{Case 1: $\mu \neq 0$ (Standard Delta Method)}
Let $g(x) = x^2$. The derivative is $g'(x) = 2x$, so at the point of interest, $g'(\mu) = 2\mu$.
Applying the Delta Method:
\[ \sqrt{n}\Big((\overline{X}_n)^2-\mu^2\Big) \ \stackrel{d}{\longrightarrow} \ \mathcal{N}(0, \underbrace{(2\mu)^2 \sigma^2}_{\text{Variance}}). \]
\textbf{Interpretation:} The limiting variance $4\mu^2\sigma^2$ comes from two sources:
\begin{enumerate}
    \item $\sigma^2$: The inherent error in estimating the mean itself.
    \item $(2\mu)^2$: The ``sensitivity" of the squaring function. Small errors in $\overline{X}_n$ are amplified by the slope $2\mu$.
\end{enumerate}

\paragraph{Case 2: $\mu = 0$ (Singularity / Second-Order Delta Method)}
If $\mu=0$, the derivative $g'(\mu) = 2(0) = 0$. The Delta Method gives:
\[ \sqrt{n}((\overline{X}_n)^2 - 0) \xrightarrow{d} 0 \cdot Z = 0. \]
This means the estimator is ``super-accurate"—the errors vanish faster than $1/\sqrt{n}$. This happens because the function $x^2$ is very flat near 0; small deviations in estimation result in quadratically smaller deviations in the output.

To get a non-trivial distribution, we must scale more aggressively. Since $\sqrt{n}$ yields 0, we try scaling by $n$:
\[ n(\overline{X}_n)^2 = \left( \sqrt{n}(\overline{X}_n - 0) \right)^2. \]
From the CLT, we know $\sqrt{n}\overline{X}_n \xrightarrow{d} \mathcal{N}(0, \sigma^2)$. We can rewrite this as $\sigma Z$, where $Z \sim \mathcal{N}(0,1)$. By the Continuous Mapping Theorem:
\[ n(\overline{X}_n)^2 \xrightarrow{d} (\sigma Z)^2 = \sigma^2 Z^2 \sim \sigma^2 \chi^2_1. \]
\textbf{Key Takeaway:} At a point where $g'(b)=0$ (a singularity), the linear approximation fails. We essentially rely on a quadratic approximation, requiring a scaling of $a_n^2$ (here $n$) instead of $a_n$ (here $\sqrt{n}$), often resulting in a Chi--square type distribution.
\end{ex}

\begin{ex}[Reciprocal of the Sample Mean]
Assume $\mu \neq 0$. Consider the estimator $g(\overline{X}_n) = \frac{1}{\overline{X}_n}$ for the quantity $1/\mu$.
The derivative is $g'(x) = -1/x^2$. Applying the Delta Method with $g'(\mu) = -1/\mu^2$:
\[ \sqrt{n}\left(\frac{1}{\overline{X}_n} - \frac{1}{\mu}\right) \stackrel{d}{\longrightarrow} \mathcal{N}\left(0, \left(-\frac{1}{\mu^2}\right)^2 \sigma^2\right) = \mathcal{N}\left(0, \frac{\sigma^2}{\mu^4}\right). \]

\textbf{Qualitative Check:}
\begin{itemize}
    \item If $\sigma^2$ is small, the input is accurate, so the output variance is small.
    \item If $\mu$ is close to 0, $1/x$ is very steep (high derivative). Small errors in $\overline{X}_n$ lead to massive errors in $1/\overline{X}_n$. This matches the $\frac{1}{\mu^4}$ term in the variance, which blows up as $\mu \to 0$.
\end{itemize}

\begin{note}
This result holds purely via asymptotic error propagation. It does not require the expectation $\mathsf{E}[1/\overline{X}_n]$ to exist (indeed, for normal variables, the inverse mean has no first moment). This is an ``apples and oranges" comparison: the Delta Method describes the \textit{limiting distribution} of the estimator, not the convergence of its moments.
\end{note}
\end{ex}

\section{Stochastic Landau Symbols}

Just as Landau symbols ($o(1)$, $O(1)$) are essential in deterministic analysis for handling error terms and convergence rates, we define their stochastic counterparts to streamline asymptotic arguments in statistics. The subscript $p$ denotes that the property holds ``in probability."

\begin{definition}[Little-o in probability]
We write $X_n = o_p(1)$ if the sequence converges to zero in probability:
\[ X_n \xrightarrow{p} 0. \]
\textit{Intuition:} This is the stochastic equivalent of a sequence of numbers converging to zero.
\end{definition}

\begin{definition}[Big-O in probability]
We write $X_n = O_p(1)$ if the sequence is \textbf{bounded in probability} (also called \textbf{uniformly tight}). Formally:

\[
    \lim\limits_{M \to \infty} \limsup\limits_{n \to \infty} \mathsf{P}(\|X_{n}\|\ge M) \longrightarrow 0 
\]

Equivalently, 
\[\forall \eps > 0 \: \exists M_{\eps}>0 \: : \: \sup_{n \in \mathbb{N}} \mathsf{P}(\|X_n\| > M_\epsilon) < \epsilon. \]
\end{definition}



For a single random variable, we know that probability mass is tight—we can always find a bounds $[-M, M]$ that contain $99\%$ of the mass.
For a \textit{sequence} $X_n$, ``boundedness in probability" means we can find a \textbf{single} $M$ that works for \textbf{all} $n$ simultaneously.
\begin{itemize}
    \item If the sequence ``drifts off" to infinity (mass escapes), it is not bounded.
    \item \textbf{Key Result:} If $X_n \xrightarrow{d} X$, then $X_n = O_p(1)$. Just as a convergent sequence of numbers is bounded, a sequence converging in distribution is tight.
\end{itemize}

\subsubsection{General Orders}

We extend these definitions to compare a random sequence $X_n$ against a deterministic rate sequence $R_n$ (e.g., $R_n = n^{-1/2}$).

\begin{itemize}
    \item $X_n = o_p(R_n)$ means $X_n = Y_n R_n$ where $Y_n = o_p(1)$.
    \item $X_n = O_p(R_n)$ means $X_n = Y_n R_n$ where $Y_n = O_p(1)$.
\end{itemize}
\textit{Informally:} $X_n = o_p(R_n)$ means $X_n / R_n \xrightarrow{p} 0$.

\subsubsection{Calculus of Stochastic Symbols}

These symbols allow us to manipulate ``remainders" in proofs algebraically. These equations should be read directionally (from left to right): ``A sequence with the property on the left also has the property on the right."

\begin{enumerate}[label=\roman*)]
    \item $o_p(1) \implies O_p(1)$:

    If a sequence goes to zero in probability, it is necessarily bounded in probability. The converse is false.

    \item $o_p(1) + o_p(1) = o_p(1)$: 

     The sum of two sequences converging to zero also converges to zero.

    \item $o_p(1) \cdot O_p(1) = o_p(1)$

         If you multiply something bounded ($O_p$) by something shrinking to zero ($o_p$), the result shrinks to zero. (This is Slutsky's theorem in disguise).

    \item $O_p(1) + O_p(1) = O_p(1)$

         The sum of two bounded sequences is bounded.
\end{enumerate}

\subsection*{Recommended Literature}
For further details on these asymptotic tools:
\begin{itemize}
    \item \textbf{Ferguson (1996):} \textit{A Course in Large Sample Theory}. (The lecturer's ``favorite" for this topic).
    \item \textbf{van der Vaart (1998):} \textit{Asymptotic Statistics}. (More general and advanced).
    \item \textbf{Billingsley:} For probability refreshers and convergence results (e.g., Skorokhod).
\end{itemize}
